\section{Loss whitening and torque-aligned physical coordinates}
\label{sec:whitening}

\subsection{First equalize the price of distance}

Before matrices, consider \(4x^2+y^2=1\). The ellipse is not evidence that
the \(x\) direction is geometrically shorter; one unit of \(x\) simply costs
four times as much. Coordinates \(X=2x\), \(Y=y\) turn it into
\(X^2+Y^2=1\). The same operation in three dimensions is
\begin{equation}
 \gls{sym:y}=\gls{sym:rmat}^{1/2}\gls{sym:i},\qquad
 \gls{sym:i}=\gls{sym:rmat}^{-1/2}\gls{sym:y},
 \label{eq:whitening-map}
\end{equation}
so that
\begin{equation}
 \gls{sym:pcu}=\norm{\gls{sym:y}}^2.
 \label{eq:whitened-loss}
\end{equation}
One unit of distance now costs one unit of copper loss in every direction.
Whitening preserves dimension, rank, intersections, and the value of every
physical quadratic form after congruence transformation. It hides amperes and
must therefore be followed by the inverse map before commanding currents.

\begin{figure}[tb]
\centering
\begin{tikzpicture}[scale=.95]
 \begin{scope}[xshift=-3.3cm]
  \draw[->] (-2.3,0)--(2.4,0) node[right]{$i_d$};
  \draw[->] (0,-1.55)--(0,1.65) node[above]{$i_{\mathrm{exc}}$};
  \draw[loss curve] (0,0) ellipse (1.85 and .82);
  \node at (0,-1.95){amperes: unequal loss price};
 \end{scope}
 \draw[-{Latex[length=3mm]},very thick] (-.8,0)--(.8,0)
 node[midway,above]{$\matr R^{1/2}$};
 \begin{scope}[xshift=3.3cm]
  \draw[->] (-1.65,0)--(1.75,0) node[right]{$y_d$};
  \draw[->] (0,-1.55)--(0,1.65) node[above]{$y_e$};
  \draw[loss curve] (0,0) circle (1.15);
  \node at (0,-1.95){distance: equal loss price};
 \end{scope}
\end{tikzpicture}
\caption{Loss whitening changes units, not physics. The ellipsoid becomes a
sphere because one coordinate unit is chosen to cost one unit of copper loss
in every direction.}
\label{fig:whitening}
\end{figure}


Put \(r_s=3R_s/2\), \(r_e=R_{\mathrm{exc}}\), and define the loss-weighted
active-flux coefficient vector
\begin{equation}
 \vect a=\begin{bmatrix}\Delta L/\sqrt{r_s}\\L_{\mathrm{exc}}/\sqrt{r_e}\end{bmatrix},
 \qquad \gls{sym:beta}=\norm{\vect a}.
 \label{eq:active-vector}
\end{equation}
The combination already printed in the torque equation suggests
\begin{equation}
 \gls{sym:spsi}=\gls{sym:dl}\gls{sym:id}
 +\gls{sym:lexc}\gls{sym:iexc}.
 \label{eq:active-flux}
\end{equation}
This is an active-flux or torque-producing-flux coordinate in the sense that
torque is \(k_Ti_qs_\psi\). Active flux is established synchronous-machine
language \cite{boldea2008}. Our additional step is to complete it orthogonally
in the loss metric so the unused coordinate has an exact resource meaning.

\subsection{Rotate after whitening, not before}

Let \(y_d,y_q,y_e\) denote the components of \(\gls{sym:y}\). Define
\begin{equation}
 \begin{bmatrix}\gls{sym:zpsi}\\\gls{sym:zq}\\\gls{sym:znull}\end{bmatrix}
 =\begin{bmatrix}
 a_d/\beta&0&a_e/\beta\\
 0&1&0\\
 -a_e/\beta&0&a_d/\beta
 \end{bmatrix}
 \begin{bmatrix}y_d\\y_q\\y_e\end{bmatrix}.
 \label{eq:torque-rotation}
\end{equation}
The matrix is orthogonal. The first row points along the active-flux
coefficient; the third is its perpendicular complement; the middle leaves the
quadrature coordinate untouched. Consequently
\begin{equation}
 \gls{sym:pcu}=\gls{sym:zpsi}^2+\gls{sym:zq}^2+\gls{sym:znull}^2,
 \qquad
 \gls{sym:torque}=\gls{sym:kappa}\gls{sym:zpsi}\gls{sym:zq},
 \label{eq:canonical-product}
\end{equation}
where
\begin{equation}
 \gls{sym:kappa}=\frac{\gls{sym:kt}\gls{sym:beta}}{\sqrt{3R_s/2}}.
 \label{eq:kappa}
\end{equation}
The inverse is explicit:
\begin{equation}
 y_d=\frac{a_dz_\psi-a_ez_0}{\beta},\quad y_q=z_q,\quad
 y_e=\frac{a_ez_\psi+a_dz_0}{\beta},\qquad
 \gls{sym:i}=\gls{sym:rmat}^{-1/2}\gls{sym:y}.
 \label{eq:torque-inverse}
\end{equation}

Without a voltage constraint, \(z_0\) contributes loss and no torque, so the
minimum has \(z_0=0\). For fixed positive torque, the arithmetic--geometric
mean inequality gives
\begin{equation}
 z_\psi^2+z_q^2\geq2z_\psi z_q=\frac{2M}{\kappa},
\end{equation}
with equality at \(z_\psi=z_q=\sqrt{M/\kappa}\). Thus
\begin{equation}
 P_{\min}^{(U\,\mathrm{free})}(M)=\frac{2|M|}{\kappa}.
 \label{eq:unconstrained-minloss}
\end{equation}
For negative torque, the factors have opposite signs and equal magnitude.
Voltage can make a nonzero \(z_0\) useful: although torque-neutral, it can move
the current vector toward the voltage-cylinder axis. Calling it useless
without the phrase ``when voltage is inactive'' would therefore be wrong.

\begin{figure}[tb]
\centering
\begin{tikzpicture}[scale=1.0]
 \draw[->] (-3,0)--(3.2,0) node[right]{$y_d$};
 \draw[->] (0,-2.1)--(0,2.25) node[above]{$y_e$};
 \draw[eigenvector] (0,0)--(2.15,1.25) node[above]{$z_\psi$: active flux};
 \draw[neutral axis] (0,0)--(-1.25,2.15) node[above left]{$z_0$: torque neutral};
 \draw[loss curve] (0,0) circle (1.55);
 \draw[->,gray] (1.0,0) arc[start angle=0,end angle=30,radius=1.0];
 \node[annotation] at (1.55,.25){loss-orthogonal rotation};
\end{tikzpicture}
\caption{After whitening, the active-flux coefficient selects one axis and
orthogonality uniquely supplies a torque-neutral complement. The quadrature
axis, perpendicular to this plane in 3D, remains unchanged.}
\label{fig:torque-axes}
\end{figure}


\paragraph{What became easy and what paid the price?}
Loss became a sphere and torque became one product. All parameter complexity
moved into the transformed voltage matrix
\(\bar{\matr Q}_U=\matr S\matr R^{-1/2}\matr Q_U\matr R^{-1/2}\matr S^{\mathsf T}\),
where \(\matr S\) is the rotation in \cref{eq:torque-rotation}. Three generic
symmetric forms cannot all share principal axes. Simultaneous orthogonal
diagonalization would require the whitened torque and voltage matrices to
commute, a special parameter relation rather than a structural property
\cite{horn2013}. SVD of the voltage map optimally simplifies voltage; the
torque rotation optimally simplifies torque. Neither basis wins every task.
